\section{Conclusion}

We presented \method{}, a pragmatic sentiment analyser for Vietnamese
social media that targets implicit sentiment, sarcasm, irony,
figurative language, code-switching, and mocking. \method{} extends
$8.4$k adjudicated comments with six pragmatic labels on top of the
ordinary 3-way polarity and 7-way emotion labels used in
UIT-VSFC \citep{vannguyen2018uitvsfc}, UIT-VSMEC \citep{ho2020uitvsmec},
and AIVIVN \citep{aivivn2019}; trains shared encoders (PhoBERT
\citep{phobert2020}, XLM-R \citep{conneau2020xlmr}, Vistral-7B
\citep{vistral2024}) and small instruction LMs under a multi-task loss
with homoscedastic-uncertainty weighting
\citep{kendall2018uncertainty}; and adds an auxiliary chain-of-thought
rationale head \citep{wei2022cot,hsieh2023distillstep} that is used
during training and dropped at inference.

On the new \method{} test set, \method{} improves macro-F1 by $+5.5$
over a strong PhoBERT fine-tune, with $+7.4$ on sarcasm and $+5.9$ on
implicit sentiment, and beats a few-shot GPT-4o-mini prompt by
$+7.2$ on sarcasm and $+3.2$ on implicit sentiment. Ordinary
sentiment is retained on all three public datasets and improves by
$+2.4$ on UIT-VSMEC. Expected calibration error is halved.
Multi-task training and the rationale auxiliary each contribute
about half of the gain; together they make the difference between a
deployable Vietnamese pragmatic classifier and a system that
collapses on sarcastic surface positives.

We release the new pragmatic annotation, the trained models, the
result JSON files, the figure-generation script, and a claim ledger
that maps every numeric claim in the paper to a source file, so that
every result can be reproduced from the released artifacts.
